\chapter{왜 VLA인가: Vision, Language, Action의 결합}

\section{장의 목표}

Vision-Language-Action(VLA) 모델은 비전 입력과 자연어 지시를 받아 로봇의 행동을 직접 생성하려는 정책 모델이다. 이 장의 목표는 VLA를 ``큰 VLM에 로봇 action head를 붙인 모델'' 정도로 축소하지 않고, 공학적으로 어떤 문제를 풀기 위해 등장했는지 정리하는 것이다. 특히 2024--2026년 연구 흐름에서 중요한 변화는 모델 크기 자체보다 행동 표현(action representation), 후처리 학습(post-training), 실시간 추론, 벤치마크, 실패 모드 분석으로 연구의 중심이 이동했다는 점이다.

이 장을 읽은 뒤 독자는 다음 질문에 답할 수 있어야 한다.

\begin{enumerate}
  \item VLA가 VLM, language-conditioned policy, diffusion policy와 어떻게 다른가?
  \item 로봇 정책을 조건부 확률 모델로 쓰면 어떤 입력과 출력이 명확해지는가?
  \item 왜 2024--2026년 VLA 연구에서 fine-tuning, action tokenization, efficiency가 동시에 중요해졌는가?
  \item VLA 논문을 읽을 때 success rate 하나만 보면 왜 충분하지 않은가?
\end{enumerate}

\section{VLA를 하나의 정책으로 보는 관점}

로봇 제어 문제에서 정책은 관측과 목표를 받아 행동을 내는 함수이다. 전통적인 강화학습 표기에서는 정책을 $\pi(\act_t \mid s_t)$로 쓴다. 그러나 실제 로봇에서 상태 $s_t$는 완전히 관측되지 않는다. 카메라 영상, force/torque, joint state, gripper state, 이전 행동, 언어 명령이 뒤섞여 들어온다. VLA 모델은 이 입력을 다음과 같이 정리한다.

\begin{equation}
  \pi_\theta(\act_t \mid \obs_{\le t}, \hist_{\le t}, \lang),
  \label{eq:vla_policy}
\end{equation}

여기서 $\obs_{\le t}$는 현재까지의 시각 관측, $\hist_{\le t}$는 proprioception과 이전 행동을 포함한 로봇 내부 이력, $\lang$은 자연어 지시, $\act_t$는 현재 시점의 로봇 행동이다. 이 표기는 VLA가 단순한 멀티모달 인식 모델이 아니라 \emph{언어 조건부 폐루프 정책}임을 드러낸다. VLM은 이미지와 텍스트의 의미를 맞추는 데 강하지만, 로봇 정책은 그 의미를 시간적으로 안정적인 행동으로 바꾸어야 한다.

\begin{conceptbox}{핵심 정의}
이 책에서 VLA 모델은 시각 관측 $\obs$, 언어 지시 $\lang$, 로봇 이력 $\hist$를 조건으로 하여 행동 $\act$ 또는 행동 시퀀스를 생성하는 학습 기반 정책 $\pi_\theta$를 의미한다. 출력이 실제 로봇 제어 입력으로 연결되지 않으면, 그것은 VLA라기보다 VLM 기반 planner 또는 captioning system에 가깝다.
\end{conceptbox}

VLA를 정책으로 보면 논문 평가 기준도 달라진다. 모델이 지시를 잘 설명하는지보다, 지시가 주어진 환경에서 행동을 얼마나 안정적으로 실행하는지가 중요하다. 같은 명령 ``빨간 컵을 잡아라''에 대해 VLM은 컵의 위치를 설명할 수 있으면 충분할 수 있다. VLA는 grasp pose, arm trajectory, gripper timing, collision avoidance, 실패 후 recovery까지 포함한 행동을 생성해야 한다. 따라서 VLA의 성능은 언어 이해 성능만으로 환원되지 않는다.

\section{왜 VLM만으로는 부족한가}

VLM은 이미지와 언어의 alignment를 학습한다. 대표적인 입력-출력 구조는 다음 조건부 분포로 나타낼 수 있다.

\begin{equation}
  p_\theta(y \mid I, x),
\end{equation}

여기서 $I$는 이미지, $x$는 텍스트 질문 또는 지시, $y$는 텍스트 응답이다. 이 구조는 시각 질의응답, captioning, grounding에는 강하지만 로봇 제어에는 세 가지 결손이 있다.

첫째, 출력 공간이 행동 공간이 아니다. 로봇은 텍스트 문장을 실행하지 않는다. 실제 실행에는 joint velocity, end-effector pose, gripper command, waypoint, torque 등 구체적 제어 변수가 필요하다. 둘째, 시간 축이 충분히 모델링되지 않는다. 한 장의 이미지에서 옳은 답을 말하는 것과, 5초 동안 접촉 상태가 바뀌는 물체를 안정적으로 조작하는 것은 다른 문제이다. 셋째, 실패의 비용이 다르다. caption이 조금 틀리는 것과 로봇이 물체를 떨어뜨리거나 충돌하는 것은 공학적으로 다른 실패이다.

따라서 VLA는 VLM의 semantic prior를 활용하되, 그 위에 행동 생성과 제어 안정성의 요구를 추가해야 한다. 이 차이를 무시하면 VLA 연구가 ``큰 모델을 붙였더니 잘된다''는 설명으로 흐른다. 수업용 관점에서는 다음 네 가지를 항상 분리해서 봐야 한다.

\begin{enumerate}
  \item \textbf{Perception}: 장면에서 필요한 객체, 관계, affordance를 찾는가?
  \item \textbf{Language grounding}: 자연어 지시가 어떤 목표 상태와 제약을 의미하는가?
  \item \textbf{Action generation}: 목표를 어떤 행동 표현으로 변환하는가?
  \item \textbf{Closed-loop correction}: 실행 중 관측이 바뀔 때 행동을 수정하는가?
\end{enumerate}

VLM은 앞의 두 항목에 강하고, VLA는 네 항목을 하나의 정책으로 묶으려 한다. 이 묶음이 성공하려면 모델 구조뿐 아니라 데이터셋, action representation, inference frequency, robot embodiment가 함께 맞아야 한다.

\section{VLA 이전의 세 흐름}

VLA는 갑자기 등장한 분야가 아니다. 세 흐름이 합쳐진 결과로 보는 것이 정확하다.

\subsection{Language-conditioned robot policy}

첫 번째 흐름은 언어 조건부 로봇 정책이다. 이 계열은 자연어 명령을 task embedding으로 바꾸고, 로봇 정책이 이를 조건으로 행동하도록 학습한다. 장점은 실제 로봇 과제에 가까운 formulation을 갖는다는 점이다. 약점은 언어와 시각의 일반화가 제한적이고, 데이터셋이 특정 로봇이나 환경에 묶이기 쉽다는 점이다.

이 흐름의 기본 학습 목적은 behavior cloning으로 쓸 수 있다.

\begin{equation}
  \mathcal{L}_{\mathrm{BC}}(\theta)
  =
  -\sum_{i=1}^{N}\sum_{t=1}^{T_i}
  \log \pi_\theta(\act^{(i)}_t \mid \obs^{(i)}_{\le t}, \hist^{(i)}_{\le t}, \lang^{(i)}).
  \label{eq:bc}
\end{equation}

이 식은 간단하지만 중요한 사실을 보여 준다. VLA는 결국 demonstration data의 분포를 따라 행동을 학습한다. 따라서 데이터가 부족하거나 특정 embodiment에 치우치면 모델이 아무리 커도 일반화가 제한된다.

\subsection{Vision-language model}

두 번째 흐름은 VLM이다. VLM은 대규모 이미지-텍스트 데이터로부터 semantic prior를 얻는다. VLA 연구가 VLM을 끌어온 이유는 로봇 데이터만으로는 세상의 객체, 도구, 관계, 언어 표현을 충분히 학습하기 어렵기 때문이다. 예를 들어 ``투명한 컵'', ``전자레인지 옆의 그릇'', ``오른쪽 손잡이가 있는 서랍'' 같은 지시는 로봇 데이터셋보다 웹 규모의 이미지-텍스트 데이터에서 더 풍부하게 나타난다.

그러나 VLM 지식은 곧바로 제어 지식이 되지 않는다. 컵이 무엇인지 아는 것과 컵을 미끄러뜨리지 않고 잡는 것은 다르다. 이 간극이 VLA의 중심 문제이다.

\subsection{Diffusion policy와 trajectory modeling}

세 번째 흐름은 diffusion policy, action chunking, trajectory modeling이다. 로봇 행동은 단일 token보다 연속 궤적에 가깝다. 특히 조작 과제에서는 접촉이 발생하고, 작은 제어 오차가 다음 단계로 누적된다. 이 때문에 행동을 한 step씩 autoregressive하게 내는 방식만으로는 충분하지 않을 수 있다.

행동 chunk를 쓰면 정책은 다음과 같이 $H$ step 행동 시퀀스를 한 번에 생성한다.

\begin{equation}
  \Act_{t:t+H-1}
  =
  (\act_t, \act_{t+1}, \ldots, \act_{t+H-1})
  \sim
  \pi_\theta(\cdot \mid \obs_{\le t}, \hist_{\le t}, \lang).
  \label{eq:action_chunk}
\end{equation}

이 방식은 inference call 수를 줄이고 short-horizon consistency를 높일 수 있다. 반대로 환경 변화에 즉각 반응하기 어렵고, chunk 안의 오류가 누적될 수 있다. 따라서 VLA에서 action chunking은 단순 속도 최적화가 아니라 안정성과 반응성 사이의 trade-off 문제이다.

\begin{sourcebox}{Pre-VLA와 VLA의 구분}
SayCan, Diffusion Policy, ACT/ALOHA는 이 책에서 VLA 자체라기보다 선행 흐름으로 분류한다. SayCan은 language-to-skill selection과 affordance 결합을, Diffusion Policy는 continuous action distribution 학습을, ACT/ALOHA는 action chunking과 bimanual imitation learning을 보여 주는 antecedent이다. 반면 RT-2, Octo, OpenVLA, \(\pi_0\), OpenVLA-OFT, SmolVLA, GR00T N1, Gemini Robotics 계열은 vision-language representation이 robot action interface로 직접 연결되는 VLA 또는 robot foundation policy로 읽는다.
\end{sourcebox}

\section{2024--2026년에 무엇이 바뀌었는가}

2024--2026년 VLA 연구의 변화는 세 문장으로 요약할 수 있다. 첫째, 연구자들은 VLM의 semantic prior를 로봇 정책에 이식하기 시작했다. 둘째, 단순히 큰 모델을 쓰는 것만으로는 success rate, latency, robustness가 충분하지 않다는 점이 드러났다. 셋째, 연구의 중심이 model scaling에서 action representation, fine-tuning protocol, benchmark, safety로 이동했다.

수집한 논문군을 보면 이 변화가 제목 수준에서도 확인된다. action, policy, diffusion, flow, trajectory와 관련된 논문이 가장 큰 묶음을 형성하고, manipulation 논문이 그다음으로 크다. efficiency 관련 논문도 매우 두껍다. 이는 VLA가 이제 ``무엇을 이해하는가''보다 ``어떻게 행동을 내는가''를 중심으로 재편되고 있음을 의미한다.

\begin{table}[t]
  \centering
  \caption{2024--2026년 VLA 연구를 읽기 위한 주요 축}
  \label{tab:vla_axes}
  \begin{tabularx}{\textwidth}{>{\bfseries}p{33mm}X}
    \toprule
    축 & 핵심 질문 \\
    \midrule
    행동 표현 & 연속 행동을 token, chunk, trajectory, diffusion sample 중 무엇으로 표현하는가? \\
    후처리 학습 & pretrained VLM/VLA를 특정 로봇 과제에 어떻게 fine-tuning 또는 post-training하는가? \\
    효율성 & 로봇 제어 주기를 만족할 만큼 inference latency와 계산량을 줄였는가? \\
    조작 성능 & contact-rich, bimanual, dexterous task에서 안정적으로 성공하는가? \\
    벤치마크 & LIBERO, Open X-Embodiment, RoboTwin 등에서 공정하게 비교되는가? \\
    안전성 & 분포 변화, 언어 모호성, adversarial instruction, 장기 horizon 실패를 다루는가? \\
    \bottomrule
  \end{tabularx}
\end{table}

이 표의 축은 장식적 분류가 아니다. VLA 논문을 읽을 때 contribution이 어느 축을 건드리는지 바로 표시할 수 있어야 한다. 예를 들어 VLA-Adapter류 논문은 후처리 학습과 효율성에 걸쳐 있고, VQ-VLA류 논문은 행동 표현에 집중한다. RoboTwin류 논문은 모델보다 벤치마크와 데이터 생성이 핵심이다. MemoryVLA나 CogVLA류 논문은 reasoning/planning 축으로 읽어야 한다.

\section{성능 지표를 어떻게 볼 것인가}

VLA 논문에서 가장 자주 보이는 지표는 task success rate이다. $M$개의 episode 중 성공한 episode 수를 $S$라고 하면 성공률은 다음과 같다.

\begin{equation}
  \mathrm{SR} = \frac{S}{M}.
  \label{eq:success_rate}
\end{equation}

그러나 성공률만으로는 좋은 VLA인지 판단하기 어렵다. 같은 성공률이라도 episode 길이, 환경 다양성, 초기 상태 분포, object set, language template, robot platform이 다르면 난이도가 달라진다. 또한 VLA는 latency가 실제 제어 성능에 직접 영향을 준다. 정책 호출 시간이 $\tau_{\mathrm{model}}$이고 로봇 제어 주기가 $\tau_{\mathrm{ctrl}}$일 때, 단순화하면 폐루프 제어가 안정적으로 동작하려면 다음 조건이 필요하다.

\begin{equation}
  \tau_{\mathrm{model}} + \tau_{\mathrm{io}} + \tau_{\mathrm{post}} \le \tau_{\mathrm{ctrl}},
  \label{eq:latency_budget}
\end{equation}

여기서 $\tau_{\mathrm{io}}$는 센서 입력과 전처리 시간, $\tau_{\mathrm{post}}$는 action decoding과 safety filtering 시간을 의미한다. 이 부등식이 깨지면 모델이 offline benchmark에서 높게 나와도 실제 로봇에서는 늦게 반응한다. 따라서 2025--2026년에 efficient VLA, speculative decoding, action chunking, lightweight adapter가 중요한 연구 축이 된 것은 자연스러운 흐름이다.

또 하나의 지표는 일반화이다. 학습 task 집합을 $\mathcal{T}_{\mathrm{train}}$, 평가 task 집합을 $\mathcal{T}_{\mathrm{test}}$라고 하자. 좋은 VLA는 같은 task의 반복이 아니라, object, scene, instruction, embodiment가 변해도 성능이 유지되어야 한다.

\begin{equation}
  \Delta_{\mathrm{gen}}
  =
  \mathbb{E}_{\tau \sim \mathcal{T}_{\mathrm{train}}}[\mathrm{SR}(\tau)]
  -
  \mathbb{E}_{\tau \sim \mathcal{T}_{\mathrm{test}}}[\mathrm{SR}(\tau)].
  \label{eq:generalization_gap}
\end{equation}

$\Delta_{\mathrm{gen}}$가 작을수록 train-test gap이 작다. 많은 VLA 논문은 평균 성공률을 강조하지만, 실제로는 어떤 조건에서 gap이 커지는지가 더 중요하다. 예를 들어 unseen object에는 강하지만 unseen instruction에는 약할 수 있고, 시뮬레이션에서는 강하지만 real robot에서는 약할 수 있다.

\section{VLA 연구를 기술적으로 읽는 기준}

이 책에서는 VLA 논문을 다음 순서로 읽는다. 첫째, 입력과 출력의 타입을 확인한다. 둘째, 행동 표현을 확인한다. 셋째, 학습 데이터와 loss를 확인한다. 넷째, inference 절차와 latency를 확인한다. 다섯째, benchmark와 baseline의 공정성을 확인한다. 여섯째, 실패 모드와 공개 자원을 확인한다.

\begin{conceptbox}{읽기 체크리스트}
VLA 논문을 읽을 때 가장 먼저 볼 것은 모델 이름이 아니라 입력, 출력, 학습 신호, 추론 주기이다. 이 네 항목이 불명확하면 contribution이 좋아 보여도 재현성과 실제 로봇 적용 가능성을 판단하기 어렵다.
\end{conceptbox}

입력은 이미지 한 장인지, multi-view인지, 비디오인지, proprioception을 포함하는지에 따라 난이도가 달라진다. 출력은 end-effector delta pose인지, joint command인지, gripper open/close인지, waypoint인지, discretized action token인지에 따라 제어 특성이 달라진다. 학습 신호는 demonstration imitation인지, reward feedback인지, preference인지, simulator verification인지에 따라 논문의 주장 범위가 달라진다. 추론 절차는 closed-loop인지 open-loop인지, action chunk를 몇 step 쓰는지, re-planning 주기가 얼마인지 확인해야 한다.

이 기준을 적용하면 같은 VLA라는 이름을 쓰는 논문도 매우 다른 문제를 풀고 있음을 알 수 있다. 어떤 논문은 사실상 VLM 기반 high-level planner이고, 어떤 논문은 low-level visuomotor policy이다. 어떤 논문은 benchmark paper이고, 어떤 논문은 action representation paper이다. 이 구분을 하지 않으면 분야가 빠르게 커질수록 논문들이 모두 비슷해 보인다.

\section{왜 manipulation이 중심이 되었는가}

수집한 논문군에서 manipulation은 가장 두꺼운 응용 축이다. 이유는 명확하다. 조작 과제는 VLA의 장점과 약점을 동시에 드러낸다. 언어 지시가 중요하고, 시각 grounding이 필요하며, 행동은 연속적이고, 접촉과 실패가 자주 발생한다. 따라서 pick-and-place, drawer opening, bimanual manipulation, dexterous grasping, contact-rich task는 VLA 모델을 평가하기 좋은 시험대가 된다.

조작 과제에서는 단순히 목표 객체를 찾는 것만으로 성공하지 않는다. grasp affordance, approach direction, force control, collision, object dynamics가 모두 영향을 준다. 예를 들어 ``컵을 들어 올려 접시 위에 놓아라''라는 명령은 객체 인식, 컵 손잡이 위치, gripper pose, lift trajectory, place pose, release timing을 포함한다. VLA가 이 과정을 하나의 정책으로 처리하려면 language grounding과 motion generation 사이의 정보 손실을 줄여야 한다.

이 때문에 2025--2026년에는 dexterous VLA, force-aware VLA, diffusion expert, bimanual benchmark가 많이 등장한다. 이 흐름은 VLA가 단순 semantic policy에서 물리 상호작용 policy로 이동하고 있음을 보여 준다.

\section{효율성이 독립적인 연구 주제가 된 이유}

큰 VLA 모델은 offline evaluation에서는 매력적이지만, 실제 로봇에서는 계산 시간이 곧 제어 문제이다. 모델이 1초에 한 번만 행동을 낸다면 빠르게 변하는 접촉 상황에 대응하기 어렵다. 이 문제를 해결하기 위해 두 방향이 등장했다.

첫째, action chunking 또는 trajectory output을 통해 모델 호출 빈도를 줄인다. 둘째, 모델 내부를 sparsification, routing, adapter, speculative decoding으로 경량화한다. 두 방법은 서로 보완적이지만 같은 문제를 풀지는 않는다. action chunking은 제어 인터페이스를 바꾸는 방법이고, 경량화는 모델 계산을 줄이는 방법이다.

효율성 연구를 평가할 때는 다음 지표를 함께 봐야 한다.

\begin{equation}
  J(\theta)
  =
  \mathrm{SR}(\theta)
  -
  \lambda_1 \cdot \mathrm{Latency}(\theta)
  -
  \lambda_2 \cdot \mathrm{Memory}(\theta)
  -
  \lambda_3 \cdot \mathrm{FailureRisk}(\theta).
  \label{eq:utility}
\end{equation}

이 식은 실제 논문들이 반드시 쓰는 objective는 아니지만, 공학적 판단에는 유용하다. 성공률이 조금 높아도 latency가 크고 메모리가 과하면 배포 가치가 낮을 수 있다. 반대로 작은 모델이 성공률을 약간 잃더라도 real-time control을 만족하면 실제 시스템에서는 더 유리할 수 있다.

\section{이 책의 관점}

이 책은 VLA를 낙관적으로 소개하지 않는다. VLA는 매우 유망하지만 아직 해결되지 않은 공학 문제가 많다. 특히 다음 네 가지는 이후 장에서 반복해서 다룬다.

\begin{enumerate}
  \item \textbf{행동 표현}: 로봇 행동을 LLM/VLM이 다룰 수 있는 형식으로 바꾸는 과정에서 정보가 손실될 수 있다.
  \item \textbf{데이터 분포}: 대규모 robot dataset이라도 embodiment와 task coverage가 제한되면 일반화가 흔들린다.
  \item \textbf{실시간성}: VLA가 실제 제어 주기를 만족하지 못하면 높은 benchmark score가 실제 성능으로 이어지지 않는다.
  \item \textbf{검증}: benchmark가 약하면 모델의 발전과 데이터셋 bias를 구분하기 어렵다.
\end{enumerate}

따라서 좋은 VLA 연구는 새로운 모델명을 제안하는 것만으로 충분하지 않다. 어떤 행동 표현을 선택했는지, 왜 그 표현이 더 안정적인지, 어떤 데이터와 baseline으로 검증했는지, latency와 실패 모드를 어떻게 보고했는지 설명해야 한다.

\section{데이터 관점에서 본 VLA}

VLA의 성능은 모델 구조만으로 결정되지 않는다. 실제로는 데이터 분포가 모델의 행동 가능 공간을 강하게 제한한다. demonstration dataset을

\begin{equation}
  \mathcal{D}
  =
  \{(\obs^{(i)}_{1:T_i}, \hist^{(i)}_{1:T_i}, \lang^{(i)}, \act^{(i)}_{1:T_i})\}_{i=1}^{N}
\end{equation}

라고 하자. 학습된 정책은 $\mathcal{D}$가 포함한 객체, 장면, 언어 표현, 로봇 kinematics, 인간 demonstrator의 습관을 모두 반영한다. 따라서 VLA 논문에서 ``large-scale robot data''라는 표현을 보더라도 다음 네 가지를 따로 확인해야 한다.

\begin{enumerate}
  \item \textbf{Embodiment coverage}: 하나의 로봇팔인지, 여러 embodiment인지, mobile manipulator나 humanoid까지 포함하는지.
  \item \textbf{Task coverage}: 단순 pick-and-place인지, tool use, contact-rich manipulation, long-horizon task까지 포함하는지.
  \item \textbf{Language coverage}: template instruction인지, 사람이 자유롭게 쓴 지시인지, ambiguity가 포함되는지.
  \item \textbf{Failure coverage}: 성공 demonstration만 있는지, 실패와 recovery trajectory도 있는지.
\end{enumerate}

특히 실패 데이터의 부재는 실제 배포에서 큰 문제가 된다. behavior cloning은 expert trajectory를 모방하는 데 강하지만, policy가 expert distribution 밖으로 벗어났을 때 복구하는 방법을 자동으로 배우지는 못한다. 이를 distribution shift로 쓰면, 학습 데이터 분포 $p_{\mathcal{D}}(\obs,\hist,\lang)$와 배포 환경 분포 $p_{\mathrm{deploy}}(\obs,\hist,\lang)$ 사이의 차이가 커질수록 성능이 떨어진다.

\begin{equation}
  \epsilon_{\mathrm{shift}}
  =
  d\!\left(
  p_{\mathcal{D}}(\obs,\hist,\lang),
  p_{\mathrm{deploy}}(\obs,\hist,\lang)
  \right),
\end{equation}

여기서 $d(\cdot,\cdot)$는 KL divergence, Wasserstein distance, 또는 실험적으로 정의한 domain gap metric일 수 있다. 중요한 것은 어떤 거리 함수를 쓰느냐보다, 논문이 train distribution과 deployment distribution의 차이를 얼마나 정직하게 다루는가이다.

\section{Pretraining, Fine-tuning, Post-training의 구분}

VLA 연구를 읽을 때 pretraining, fine-tuning, post-training을 섞어 쓰면 논문 주장이 흐려진다. 이 책에서는 다음처럼 구분한다.

\begin{table}[t]
  \centering
  \caption{VLA 학습 단계의 구분}
  \label{tab:training_stages}
  \begin{tabularx}{\textwidth}{>{\bfseries}p{31mm}X}
    \toprule
    단계 & 공학적 의미 \\
    \midrule
    Pretraining & 웹 이미지-텍스트, 로봇 trajectory, 멀티모달 데이터에서 일반 표현과 기본 행동 prior를 얻는 단계이다. \\
    Fine-tuning & 특정 robot embodiment, action space, task family에 맞추어 파라미터 전체 또는 일부를 조정하는 단계이다. \\
    Adapter tuning & backbone은 대부분 고정하고 작은 모듈만 학습하여 비용과 catastrophic forgetting을 줄이는 단계이다. \\
    Post-training & instruction following, reward feedback, simulator verification, human preference 등을 이용해 실행 품질을 보정하는 단계이다. \\
    Deployment calibration & 실제 로봇 주기, safety filter, controller interface, sensor delay에 맞추어 policy를 시스템에 연결하는 단계이다. \\
    \bottomrule
  \end{tabularx}
\end{table}

OpenVLA-OFT류의 문제의식은 이 구분에서 fine-tuning과 deployment calibration 사이에 위치한다. 즉, pretrained VLA가 이미 의미 있는 prior를 갖고 있더라도, 실제 task success와 inference speed를 동시에 만족하려면 action head, training recipe, inference path를 다시 설계해야 한다. 이 관점에서는 VLA의 발전을 단순히 parameter scaling으로 설명할 수 없다.

학습 목적도 하나로 고정되지 않는다. behavior cloning loss에 action regularization이나 temporal smoothness를 추가할 수 있다.

\begin{equation}
  \mathcal{L}(\theta)
  =
  \mathcal{L}_{\mathrm{BC}}(\theta)
  +
  \alpha \sum_{t=2}^{T}
  \lVert \hat{\act}_{t} - \hat{\act}_{t-1} \rVert_2^2
  +
  \beta \mathcal{L}_{\mathrm{aux}}(\theta).
  \label{eq:vla_loss}
\end{equation}

두 번째 항은 행동이 시점마다 불안정하게 흔들리는 것을 줄이는 smoothness term이고, $\mathcal{L}_{\mathrm{aux}}$는 language grounding, object detection, affordance prediction, value estimation 같은 보조 손실일 수 있다. 논문을 읽을 때는 정확히 어떤 loss가 success rate 향상에 기여했는지 ablation이 있는지 확인해야 한다.

\section{실패 모드의 분해}

VLA 실패는 하나의 숫자로 요약하기 어렵다. 실패 원인을 분해하지 않으면 모델 개선 방향도 모호해진다. 예를 들어 success rate가 60\%에서 70\%로 올랐다고 해도, 실패가 perception에서 줄었는지, grasp planning에서 줄었는지, language ambiguity에서 줄었는지 알 수 없다.

\begin{table}[t]
  \centering
  \caption{VLA 실패 모드의 기술적 분해}
  \label{tab:failure_modes}
  \begin{tabularx}{\textwidth}{>{\bfseries}p{32mm}X}
    \toprule
    실패 위치 & 확인해야 할 증상 \\
    \midrule
    Perception failure & 목표 객체를 찾지 못하거나 distractor를 목표로 잘못 선택한다. \\
    Grounding failure & 언어 지시의 관계, 순서, 제약 조건을 잘못 해석한다. \\
    Action representation failure & 의미는 맞지만 action token, waypoint, trajectory가 실제 제어에 부적절하다. \\
    Dynamics failure & 접촉, 마찰, 무게, deformable object 때문에 계획한 행동이 물리적으로 실패한다. \\
    Temporal failure & 초기 step은 맞지만 long-horizon에서 오차가 누적되어 task가 깨진다. \\
    Recovery failure & 실패 후 관측을 보고도 재시도, re-grasp, re-plan을 하지 못한다. \\
    Safety failure & 충돌, 과도한 force, 금지된 객체 조작 등 시스템 제약을 위반한다. \\
    \bottomrule
  \end{tabularx}
\end{table}

이 분해는 논문 리뷰에도 직접 쓸 수 있다. 강한 VLA 논문은 평균 성능뿐 아니라 실패 위치를 보여 준다. 약한 논문은 성공률 표만 제시하고, 왜 실패했는지 거의 말하지 않는다. 특히 실제 로봇 배포를 주장하는 논문이라면 safety failure와 recovery failure를 따로 보고해야 한다.

\section{VLA 논문의 최소 재현 단위}

공학 수업에서 VLA 논문을 다룰 때는 결과를 믿기 전에 최소 재현 단위를 확인해야 한다. 최소 재현 단위는 다음 여섯 항목이다.

\begin{enumerate}
  \item 모델 backbone과 학습 가능한 모듈의 범위.
  \item 관측 modality와 카메라 구성.
  \item action space의 정의와 controller interface.
  \item 학습 데이터의 출처, task split, train-test split.
  \item evaluation protocol과 success 판정 기준.
  \item 코드, checkpoint, dataset, benchmark script 공개 여부.
\end{enumerate}

이 여섯 항목 중 action space와 evaluation protocol이 빠지면 로봇 논문으로서 재현성이 크게 떨어진다. 예를 들어 end-effector delta pose를 쓰는 정책과 joint velocity를 쓰는 정책은 같은 success rate라도 제어 난이도가 다르다. 사람이 reset한 episode와 자동 reset된 benchmark도 비교가 어렵다. 따라서 이 책에서는 각 대표 논문을 소개할 때 모델 그림보다 먼저 이 재현 단위를 표로 정리할 것이다.

\section{작업 예제: 컵을 들어 접시에 놓기}

VLA의 차이를 가장 단순하게 보려면 하나의 조작 과제를 수식과 시스템 구성으로 분해하면 된다. 명령은 ``빨간 컵을 들어 접시 위에 놓아라''라고 하자. 환경에는 빨간 컵, 파란 컵, 접시, 장애물이 있고, 로봇은 wrist camera와 external camera를 사용한다. 이 과제는 겉으로는 단순하지만 VLA의 주요 병목을 모두 포함한다.

먼저 VLM planner 방식은 다음처럼 동작할 수 있다.

\begin{equation}
  z_{1:K}
  =
  f_{\mathrm{VLM}}(I_t, \lang),
\end{equation}

여기서 $z_{1:K}$는 ``빨간 컵 찾기'', ``컵 쪽으로 이동'', ``잡기'', ``접시 위로 이동'', ``놓기'' 같은 symbolic subgoal이다. 이 방식은 설명 가능하고 high-level planning에는 편하지만, 각 subgoal을 실제 제어 명령으로 바꾸는 별도 motion planner와 controller가 필요하다.

반대로 low-level visuomotor policy는 symbolic subgoal을 거치지 않고 바로 행동을 낸다.

\begin{equation}
  \act_t
  =
  g_\theta(I_t, q_t, \lang),
\end{equation}

여기서 $q_t$는 joint state 또는 end-effector state이다. 이 방식은 제어에 직접 연결되지만, 언어 지시의 compositional structure를 잘 활용하지 못할 수 있다. VLA는 두 접근의 중간에 있다. VLM의 semantic prior를 쓰되, 출력은 실제 action space로 내려보낸다.

\begin{equation}
  \Act_{t:t+H-1}
  =
  \pi_\theta(I^{\mathrm{ext}}_t, I^{\mathrm{wrist}}_t, q_t, \lang).
\end{equation}

이 예제에서 좋은 VLA 논문이라면 다음 정보를 명확히 말해야 한다. 빨간 컵과 파란 컵을 어떻게 구분했는가? 접시 위라는 spatial relation을 어떤 표현으로 grounding했는가? grasp는 pose regression인가, action token인가, diffusion sample인가? 컵을 잡은 뒤 wrist camera의 시야가 바뀔 때 policy가 closed-loop로 수정되는가? 컵이 미끄러졌을 때 recovery를 하는가? 이 질문에 답하지 못하면 성공 영상이 있어도 공학적으로 충분한 설명이 아니다.

\begin{table}[t]
  \centering
  \caption{컵-접시 과제에서 세 접근의 차이}
  \label{tab:cup_plate_example}
  \begin{tabularx}{\textwidth}{>{\bfseries}p{30mm}XXX}
    \toprule
    항목 & VLM planner & Visuomotor policy & VLA policy \\
    \midrule
    출력 & symbolic subgoal & low-level action & language-conditioned action \\
    장점 & 해석 가능성 & 빠른 제어 연결 & 의미 이해와 행동 생성의 결합 \\
    약점 & controller 의존 & 언어 일반화 약함 & 데이터, latency, action 표현 모두 민감 \\
    평가 포인트 & plan validity & control success & task success, latency, generalization \\
    \bottomrule
  \end{tabularx}
\end{table}

이 예제는 이후 장의 기준 예제로 계속 사용할 수 있다. action tokenization 장에서는 $\Act_{t:t+H-1}$를 어떻게 이산화하는지 볼 것이고, fine-tuning 장에서는 이 과제에 맞추어 어떤 파라미터를 조정하는지 볼 것이다. efficiency 장에서는 이 policy가 제어 주기 안에 실행되는지 계산하고, robustness 장에서는 distractor, occlusion, language ambiguity가 들어왔을 때 어떤 실패가 생기는지 분석한다.

\section{수업에서 요구할 산출물}

이 책을 실제 대학원 수업에서 쓴다면 학생에게 단순 발표보다 다음 산출물을 요구하는 편이 낫다. 첫째, 선택한 논문의 policy interface를 한 장짜리 표로 정리한다. 둘째, action representation을 수식으로 쓴다. 셋째, benchmark protocol을 train/test split과 success criterion 중심으로 재작성한다. 넷째, 실패 모드를 최소 세 가지로 분해한다. 다섯째, 공개 코드와 checkpoint가 없을 경우 재현 가능한 최소 실험을 제안한다.

이 요구사항은 논문을 비판적으로 읽기 위한 최소 조건이다. VLA 논문은 모델 구조 그림이 복잡해 보일수록 핵심이 흐려질 수 있다. 반대로 입력, 출력, loss, benchmark가 명확하면 그림이 단순해도 좋은 연구일 수 있다. 공학 수업에서 중요한 것은 설명의 화려함이 아니라, 같은 조건에서 다른 연구자가 실험을 반복할 수 있는가이다.

\section{강의에서의 사용 방식}

이 장은 첫 주 강의의 기준선을 만든다. 학생들은 이후 장에서 나오는 논문을 읽을 때 매번 같은 질문으로 돌아와야 한다.

\begin{enumerate}
  \item 이 논문은 VLA의 어느 병목을 건드리는가?
  \item action representation은 무엇이고, 왜 그것을 선택했는가?
  \item success rate 외에 latency, robustness, generalization을 보고하는가?
  \item baseline은 충분히 강한가?
  \item 실제 로봇에서 실패할 가능성이 가장 큰 위치는 어디인가?
\end{enumerate}

강의 과제로는 한 논문을 골라 \cref{tab:failure_modes}의 실패 모드 표에 맞춰 다시 분석하게 할 수 있다. 이 작업은 단순 요약보다 더 중요하다. VLA 분야에서는 논문 수가 빠르게 늘기 때문에, 모델 이름을 기억하는 능력보다 contribution을 병목 축에 배치하는 능력이 연구 생산성을 결정한다.

\section{요약}

VLA는 VLM의 semantic prior와 로봇 정책 학습을 결합하려는 시도이다. 그러나 VLA의 핵심은 멀티모달 이해가 아니라 행동 생성이다. 2024--2026년 연구 흐름은 이 점을 점점 더 분명하게 보여 준다. action tokenization, diffusion action decoder, action chunking, post-training, efficient inference, robustness benchmark는 모두 같은 질문으로 모인다. ``언어로 지시받은 목표를 실제 로봇이 시간 안에, 안전하게, 반복 가능하게 수행할 수 있는가?''

다음 장에서는 VLA의 기본 구조를 더 구체적으로 분해한다. 입력 modality, backbone, policy head, action representation, loss function, inference loop를 각각 따로 보고, 이후 장에서 등장할 논문들을 비교할 공통 표기법을 만든다.

\begin{sourcebox}{Key papers}
\begin{itemize}
  \item Ahn et al. (2022), SayCan, pre-VLA language-to-skill grounding: LLM planning과 robot affordance 결합을 이해하기 위한 출발점이다.
  \item Brohan et al. (2022), RT-1, robot transformer: 대규모 robot trajectory가 language-conditioned policy 일반화에 어떤 역할을 하는지 보여 준다.
  \item Brohan et al. (2023), RT-2, VLA transfer: web-scale VLM knowledge를 robot action token으로 옮기는 대표 사례이다.
  \item Kim et al. (2024), OpenVLA, arXiv preprint: 7B open-source VLA와 970k real-world demonstration 기반 학습을 제시한 핵심 anchor이다.
  \item Black et al. (2024/2026), \(\pi_0\), RSS 2025: VLM 위에 flow matching action architecture를 얹은 general robot control anchor이다.
\end{itemize}
\end{sourcebox}

\begin{assignmentbox}{Reading assignment [Basic]}
OpenVLA와 SayCan을 같은 13-question template로 읽고, 하나는 VLA로, 다른 하나는 pre-VLA foundation으로 분류되는 이유를 controller interface와 action output 관점에서 1쪽으로 정리하라.
\end{assignmentbox}
